\documentclass[11pt,a4paper]{article}
\usepackage[left=23mm,right=23mm,top=22mm,bottom=22mm,footskip=11mm]{geometry}
\usepackage{amsmath,amssymb}
\usepackage{assets/tex/cm-bea}
\usepackage{xeCJK}
\setmainfont{Latin Modern Roman}
\setCJKmainfont{FandolSong-Regular.otf}[Path=assets/fonts/fandol/,BoldFont=FandolSong-Bold.otf,ItalicFont=FandolSong-Regular.otf]
\setCJKsansfont{FandolSong-Regular.otf}[Path=assets/fonts/fandol/,BoldFont=FandolSong-Bold.otf]
\renewcommand{\bfdefault}{bx}
\xeCJKsetup{PunctStyle=quanjiao,CheckSingle=true}
\usepackage{setspace}\setstretch{1.10}
\setlength{\parindent}{2em}\setlength{\parskip}{4pt}
\usepackage{xcolor,titlesec,tcolorbox,tikz,enumitem,longtable,booktabs,array,calc,needspace}
\tcbuselibrary{breakable,skins}\usetikzlibrary{arrows.meta,positioning}
\definecolor{sitepaper}{HTML}{FAF9F5}\pagecolor{sitepaper}
\definecolor{navy}{HTML}{10136B}\definecolor{inkgray}{HTML}{474747}\definecolor{warm}{HTML}{94662C}\definecolor{teal}{HTML}{236B66}
\titleformat{\section}{\large\bfseries\color{navy}}{}{0pt}{}
\titleformat{\subsection}{\normalsize\bfseries\color{navy}}{}{0pt}{}
\titlespacing*{\section}{0pt}{17pt}{9pt}\titlespacing*{\subsection}{0pt}{12pt}{6pt}
\newtcolorbox{sboxenv}[2]{enhanced,breakable,colframe=#1,colback=#1!3!white,colbacktitle=#1,coltitle=white,fonttitle=\bfseries,title={#2},boxrule=0.45pt,arc=1.4pt,left=8pt,right=8pt,top=6pt,bottom=6pt,before skip=9pt,after skip=9pt,pad at break*=2mm}
\setlist[enumerate]{leftmargin=2.2em,itemsep=3pt,parsep=0pt,topsep=3pt}
\setlist[itemize]{label={},leftmargin=1em,itemsep=3pt}
\providecommand{\tightlist}{\setlength{\itemsep}{3pt}\setlength{\parskip}{0pt}}
\usepackage[unicode,hidelinks,linktoc=all]{hyperref}\usepackage{bookmark}
\hypersetup{pdftitle={实变函数-v1：4.5 重积分与累次积分},pdfauthor={Vincent Arqain}}
\renewcommand{\contentsname}{本节目次}\setcounter{tocdepth}{1}\setcounter{secnumdepth}{0}
\widowpenalty=10000\clubpenalty=10000\displaywidowpenalty=10000\predisplaypenalty=5000
\allowdisplaybreaks[1]\raggedbottom\emergencystretch=2em
\begin{document}
\begin{center}{\LARGE\bfseries 实变函数-v1\par}\vspace{7pt}
{\Large\bfseries 4.5 重积分与累次积分\par}\vspace{6pt}
{第四章\quad Lebesgue 积分\par}\end{center}
沿竖线切单位正方形中的三角形 \(D=\{(x,y):0<y<x<1\}\)，截线长为
\(x\)；沿横线切，截线长为 \(1-y\)。两种切法都算出面积
\(\frac12\)。若每点还带有正负相间的密度，两种切法算出的总量还会相同吗？
\par\textbf{主线：}先把矩形的体积写成截面测度的积分，用测度正则性推广到可测集。把集合的指示函数换成非负函数，得到
Tonelli 定理。再控制正、负部，便能说明何时可以交换带符号函数的积分次序。
\begin{center}
\begin{tikzpicture}[node distance=5mm and 5mm,box/.style={draw=navy!65,fill=navy!3,rounded corners=1pt,minimum width=34mm,minimum height=10mm,align=center,font=\small},>=Stealth]
\node[box](a){截面测度};\node[box,right=of a](b){非负函数换序};\node[box,right=of b](c){绝对可积换序};\node[box,right=of c,fill=teal!5](d){分布函数与卷积};
\draw[->](a)--(b);\draw[->](b)--(c);\draw[->](c)--(d);
\end{tikzpicture}\end{center}
对应徐森林、薛春华《实变函数论》§3.7，重点为印刷第 283---291 页的定理
3.7.1---3.7.6、例 3.7.1---3.7.2，以及第 304 页的六道练习。记 \(m_k\) 为
\(\mathbb R^k\) 上的 Lebesgue 测度，\(p,q\) 为正整数。非负积分可取
\(+\infty\)，约定 \(0\cdot\infty=0\)。
\begingroup\small\setlength{\parskip}{0pt}\tableofcontents\endgroup
\newpage
\Needspace{10\baselineskip}\section[\S{}4.5.1 截面与体积]{4.5.1 截面与体积}\label{ux622aux9762ux4e0eux4f53ux79ef}

写 \(\mathbb R^{p+q}=\mathbb R^p\times\mathbb R^q\)。固定 \(x\) 或
\(y\)，集合 \(E\) 的两个截面分别为 \begin{equation*}
E_x=\{y:(x,y)\in E\},\qquad E^y=\{x:(x,y)\in E\}.
\end{equation*}
截面与补、并、交可交换，因为判断一个点是否属于截面，只需把固定坐标补回去。若
\(E\) 是 Borel 集，则每个 \(E_x\) 都是 Borel 集，它是 \(E\) 在连续映射
\(y\mapsto(x,y)\) 下的原像。对 Lebesgue
可测集，完备化引入了零集的一切子集，必须单独处理。

\begin{sboxenv}{navy}{定理 3.7.3（截面公式）}
设 \(E\subset\mathbb R^{p+q}\) 为 Lebesgue 可测集。对几乎所有
\(x\)，\(E_x\) 可测。其测度在例外零集上补零后是可测函数，且

\begin{equation*}
m_{p+q}(E)=\int_{\mathbb R^p}m_q(E_x)\,dx
=\int_{\mathbb R^q}m_p(E^y)\,dy .
\tag{4.5.1}
\end{equation*}

\end{sboxenv}

\textbf{证明}　先把 \(E\) 限在一个有限半开方体 \(Q=I\times J\)
中。对半开矩形 \(A\times B\subset Q\)，截面测度是
\(m_q(B)\mathbf1_A(x)\)，其积分是
\(m_p(A)m_q(B)\)，恰为矩形体积。对有限个互不相交的半开矩形，可将这两个等式相加，因为同一坐标下的截面也两两不交。

设 \(G\) 在全空间开。在 \(Q\) 内作逐次二分网格，以 \(G_n\) 表示第 \(n\)
级完全包含于 \(G\) 的网格小方体之并。于是
\(G_n\uparrow G\cap Q\)。事实上，每个 \(z\in G\cap Q\) 到 \(G\)
外有正的邻域半径，小方体的直径终将小于该半径。对每个
\(x\)，截面也递增。由测度的下连续性和 Levi 定理，截面公式对 \(G\cap Q\)
成立，截面测度函数可测。对与 \(Q\) 的相对补集，可从 \(m_q(J)\mathbf1_I\)
和 \(m_{p+q}(Q)\) 中分别减去这两个有限的量。因此公式也适用于
\(K\cap Q\)，其中 \(K\) 闭。

现设 \(E\subset Q\) 可测。由正则性，取紧集 \(K_n\subset E\) 和开集
\(G_n\supset E\)，使 \(m_{p+q}((G_n\cap Q)\setminus K_n)<2^{-n}\)。用前
\(n\) 个紧集的并和前 \(n\)
个开集的交替换，仍有同一估计，且内层递增、外层递减。令
\(K=\bigcup_n K_n\)、\(H=Q\cap\bigcap_n G_n\)。这两个集合的截面均为
Borel 集。内层用递增极限，外层用 \(Q\) 的有限体积和截面的统一可积控制
\(m_q(J)\mathbf1_I\)，可知 \(K,H\) 均满足公式。又
\(K\subset E\subset H\)，且 \(m_{p+q}(H\setminus K)=0\)，因而 \begin{equation*}
\int_{\mathbb R^p}\bigl[m_q(H_x)-m_q(K_x)\bigr]\,dx=0 .
\end{equation*} 被积函数非负，故对几乎所有 \(x\)，\(H_x\setminus K_x\)
为零集。由完备性，夹在两者之间的 \(E_x\) 可测且与 \(K_x\)
等测度。积分即得 (4.5.1) 的第一个等号。对无界 \(E\)，将 \(E\)
与递增至全空间的半开方体相交；并起可数个例外零集后，再用下连续性和 Levi
定理。交换坐标得到另一个等号。

\begin{sboxenv}{warm}{辨析　一个坏截面可以有多坏}
取不可测集 \(V\subset[0,1]\)。集合 \(\{0\}\times V\)
包含在平面零测线段内，因而可测，但它在 \(x=0\) 处的截面就是
\(V\)。这说明 (4.5.1) 保留``几乎所有''的原因。

\end{sboxenv}

\textbf{推论（教材定理 3.7.4）}　若
\(A\subset\mathbb R^p\)、\(B\subset\mathbb R^q\) 可测，则 \(A\times B\)
可测，且 \(m_{p+q}(A\times B)=m_p(A)m_q(B)\)。

\textbf{证明}　先说明乘积可测。若 \(N\subset\mathbb R^p\)
为零集，给定有界方体 \(J\subset\mathbb R^q\)，用总体积小于
\(\varepsilon\) 的可数个矩形覆盖 \(N\)，再与 \(J\) 作积，就以总体积小于
\(\varepsilon m_q(J)\) 的矩形覆盖了 \(N\times J\)。扩大 \(J\)，知
\(N\times\mathbb R^q\) 仍为零集。取与 \(A,B\) 分别只差零集的 Borel 集
\(A_0,B_0\)。乘积 \(A_0\times B_0\) 为 Borel 集，而它与 \(A\times B\)
的对称差包含在两个上述形式的零测柱集之并中，故 \(A\times B\)
可测。其截面在 \(x\in A\) 时为 \(B\)，其余为空集。代入 (4.5.1)
即得结论，也包括其中一个因子零测而另一个测度无穷的情形。

\Needspace{10\baselineskip}\section[\S{}4.5.2
非负函数的两种积分次序]{4.5.2
非负函数的两种积分次序}\label{ux975eux8d1fux51fdux6570ux7684ux4e24ux79cdux79efux5206ux6b21ux5e8f}

\begin{sboxenv}{navy}{定理 3.7.1（Tonelli）}
设 \(f:\mathbb R^{p+q}\to[0,\infty]\) 为 Lebesgue 可测函数。对几乎所有
\(x\)，\(f(x,\cdot)\) 可测；内积分的可测代表记为
\(F(x)\)。沿另一方向同理，且

\begin{equation*}
\begin{aligned}
\int_{\mathbb R^{p+q}}f
&=\int_{\mathbb R^p}\left(\int_{\mathbb R^q}f(x,y)\,dy\right)dx\\
&=\int_{\mathbb R^q}\left(\int_{\mathbb R^p}f(x,y)\,dx\right)dy .
\end{aligned}
\tag{4.5.2}
\end{equation*}

三个积分可以同时为 \(+\infty\)。

\end{sboxenv}

\textbf{证明}　若 \(f=\mathbf1_E\)，上述结论正是
(4.5.1)。有限次相加便得非负简单函数的结论。一般情形，取非负简单函数
\(s_n\uparrow f\)。合并所有 \(s_n\) 的截面例外集，得到同一个零集
\(N\)。对 \(x\notin N\)，截面 \(s_n(x,\cdot)\) 是可测递增列，极限
\(f(x,\cdot)\) 可测，且 \begin{equation*}
F_n(x)=\int s_n(x,y)\,dy \ \uparrow\ \int f(x,y)\,dy=F(x).
\end{equation*} 各 \(F_n\) 在 \(N\) 上统一取零，故 \(F\) 也有可测代表。外层再用一次
Levi 定理，得到
\(\int F=\lim_n\int F_n=\lim_n\int s_n=\int f\)。另一方向的证明相同。两次极限都是单调极限，所以无需预先估计积分是否有限。

\textbf{零延拓}　若只在可测集 \(D\) 上给定 \(f\ge0\)，令它在域外为零，则
(4.5.2) 的内积分只需在 \(D_x\) 或 \(D^y\) 上求。开篇的三角形便给出
\(\int_0^1 x\,dx=\int_0^1(1-y)\,dy\)。一般区域换序时，内层上下限应从同一个
\(D\) 的另一方向截面重新读出。

\Needspace{10\baselineskip}\section[\S{}4.5.3
带符号函数的换序条件]{4.5.3
带符号函数的换序条件}\label{ux5e26ux7b26ux53f7ux51fdux6570ux7684ux6362ux5e8fux6761ux4ef6}

\begin{sboxenv}{navy}{定理 3.7.2（Fubini）}
设 \(f\) 联合可测。若
\(\int_{\mathbb R^{p+q}}|f|<\infty\)，则两个方向上几乎所有的截面均绝对可积，两种内积分都是可积函数，并且
(4.5.2) 对 \(f\) 成立，三项均为有限值。

\end{sboxenv}

\textbf{证明}　对 \(|f|\) 使用 Tonelli。非负函数
\(x\mapsto\int |f(x,y)|\,dy\) 的积分有限，故它几乎处处有限。可测性可对
\(f^+\)、\(f^-\) 使用上一节的截面论证；因此几乎所有截面绝对可积。令
\(F_\pm(x)=\int f^\pm(x,y)\,dy\)。两者均可测且可积，在它们同为有限的点作差得到
\(F(x)=\int f(x,y)\,dy\)。又 \(|F(x)|\le\int |f(x,y)|\,dy\)，故 \(F\)
可积。由 (4.5.2) 及积分线性， \begin{equation*}
\int F=\int F_+-\int F_-=\int f^+-\int f^-=\int f.
\end{equation*} 横向截面给出同一结果。复值情形对实、虚部分别应用此结论。

\textbf{实用判据}　只须证明 \(|f|\) 的任意一种累次积分有限。因为 Tonelli
先对非负的 \(|f|\) 无条件给出等式，这一个有限值就等于
\(\int|f|\)。若已知某个非负可积函数控制 \(|f|\)，可直接比较。仅从
\(\int|\int f(x,y)\,dy|\,dx<\infty\)
无法作同样推断，绝对值放在内积分外面已经允许抵消。

\Needspace{10\baselineskip}\section[\S{}4.5.4
两种累次积分透露了什么]{4.5.4
两种累次积分透露了什么}\label{ux4e24ux79cdux7d2fux6b21ux79efux5206ux900fux9732ux4e86ux4ec0ux4e48}

\begin{sboxenv}{inkgray}{例 3.7.1　两个方向都算出零}
在 \([-1,1]^2\) 上取
\(f(x,y)=\frac{xy}{(x^2+y^2)^2}\)，原点取零。两种累次积分都存在且相等，\(f\)
是否可积？

\end{sboxenv}

\textbf{解}　固定一个非零坐标，截面是另一变量的连续奇函数；固定坐标为零时截面恒为零。因此两种内积分都恒为零。看第一象限，函数非负，固定
\(x>0\) 作代换 \(u=x^2+y^2\) 得 \begin{equation*}
\int_0^1\frac{xy}{(x^2+y^2)^2}\,dy
=\frac{1}{2x}-\frac{x}{2(1+x^2)} .
\end{equation*} 右端第一项在零点的积分发散，第二项有界。Tonelli
给出第一象限正部积分为无穷。关于 \(y\)
反射，第四象限的负部积分也为无穷，故原函数没有 Lebesgue
积分。这里即使两个答案已经相等，也不能倒推绝对可积。

\begin{sboxenv}{inkgray}{例　两个方向算出不同的数}
在 \(D=(0,1)^2\) 上令
\(g(x,y)=\frac{x^2-y^2}{(x^2+y^2)^2}\)。计算两种累次积分。

\end{sboxenv}

\textbf{解}　沿 \(y\) 方向，被积函数是 \(\frac{y}{x^2+y^2}\) 关于 \(y\)
的导数。沿 \(x\) 方向，它是 \(\frac{-x}{x^2+y^2}\) 关于 \(x\)
的导数。由端点代入， \begin{equation*}
\begin{aligned}
\int_0^1\left(\int_0^1 g(x,y)\,dy\right)dx
 &=\int_0^1\frac{dx}{1+x^2}=\frac{\pi}{4},\\
\int_0^1\left(\int_0^1 g(x,y)\,dx\right)dy
 &=-\int_0^1\frac{dy}{1+y^2}=-\frac{\pi}{4}.
\end{aligned}
\end{equation*} 若 \(\int_D|g|\) 有限，Fubini 将迫使两值相等。还可定位发散的区域。在
\(0<y<x/2\) 内，有
\(x^2-y^2\ge 3x^2/4\)、\((x^2+y^2)^2\le 25x^4/16\)，故
\(g\ge 12/(25x^2)\)。这条楔形在位置 \(x\) 的宽为
\(x/2\)，于是正部积分至少为
\(\int_0^1\frac{6}{25x}\,dx=\infty\)。交换两坐标使函数变号，负部同样不可积。

\Needspace{10\baselineskip}\section[\S{}4.5.5
下方图形与分布函数]{4.5.5
下方图形与分布函数}\label{ux4e0bux65b9ux56feux5f62ux4e0eux5206ux5e03ux51fdux6570}

设 \(E\subset\mathbb R^p\) 可测，\(h:E\to[0,\infty]\)。定义下方图形
\(D_h=\{(x,t):x\in E,\ 0<t<h(x)\}\)。若 \(h\) 可测，则 \begin{equation*}
D_h=\bigcup_{r\in\mathbb Q,\ r>0}\{x\in E:h(x)>r\}\times(0,r).
\end{equation*} 由乘积集可测性，\(D_h\) 可测。反过来，若 \(D_h\) 可测，其竖截面就是
\((0,h(x))\)，长度恰为 \(h(x)\)。由
(4.5.1)，截面长度有可测代表，完备性保证 \(h\) 在 \(E\)
上可测。这给出教材定理 3.7.6 的可测性判据。

\begin{sboxenv}{navy}{下方图形与分布函数公式}
令 \(\lambda_h(t)=m_p(\{x\in E:h(x)>t\})\)。沿两个方向切 \(D_h\)，得到

\begin{equation*}
m_{p+1}(D_h)=\int_E h(x)\,dx=\int_0^\infty\lambda_h(t)\,dt .
\tag{4.5.3}
\end{equation*}

\end{sboxenv}

\textbf{证明}　按 \(x\) 取截面得高度 \(h(x)\)；按 \(t>0\)
取截面得超水平集 \(\{h>t\}\)。把两种截面代入 (4.5.1) 即得。更一般地，对
\(r>0\) 把 \(h(x)^r\) 写成 \(\int_0^{h(x)}rt^{r-1}\,dt\)。在 \(t>0\)
上被积函数非负，Tonelli 给出 \begin{equation*}
\int_E h(x)^r\,dx
=r\int_0^\infty t^{r-1}\lambda_h(t)\,dt .
\tag{4.5.4}
\end{equation*} 例如 \(h(x)=x^{-1/2}\)，\(0<x<1\)。对 \(0<t\le1\) 有
\(\lambda_h(t)=1\)，对 \(t>1\) 有 \(\lambda_h(t)=t^{-2}\)。代入
(4.5.4)，第一段积分为 \(1\)，第二段为 \(r\int_1^\infty t^{r-3}\,dt\)。故
\(0<r<2\) 时结果为 \(\frac{2}{2-r}\)，\(r\ge2\)
时发散。函数在零点附近的大小，由很高的水平集缩小的速度精确记录。

\textbf{图形与下方图形（定理 3.7.5、例 3.7.2）}　若实值 \(h\)
可测，其图形 \(\Gamma_h=\{(x,t):x\in E,\ t=h(x)\}\) 可测，因为
\((x,t)\mapsto h(x)\) 在 \(E\times\mathbb R\)
上可测（水平集是可测柱集），故可与坐标函数 \(t\)
作差并取零水平集。每个竖截面至多含一个点，(4.5.1) 给出
\(m_{p+1}(\Gamma_h)=0\)。因此在 \(D_h\)
中加入上、下边界不改变体积。但只知道图形零测，无从恢复函数的可测性。取
\(V\subset[0,1]\) 不可测，\(\mathbf1_V\)
的图形包含在两条水平线段中，是平面零集；它的下方图形却不可测，否则由前面的逆推就会得到
\(V\) 可测。

\Needspace{10\baselineskip}\section[\S{}4.5.6
卷积的存在性与积分]{4.5.6
卷积的存在性与积分}\label{ux5377ux79efux7684ux5b58ux5728ux6027ux4e0eux79efux5206}

\begin{sboxenv}{navy}{性质　两个可积函数的卷积}
设 \(f,g\in L^1(\mathbb R^d)\)。对几乎所有 \(x\)，卷积
\((f*g)(x)=\int_{\mathbb R^d}f(x-y)g(y)\,dy\) 绝对收敛。取可测代表后，

\begin{equation*}
\|f*g\|_1\le\|f\|_1\|g\|_1,\qquad
\int_{\mathbb R^d}(f*g)=\left(\int_{\mathbb R^d}f\right)\left(\int_{\mathbb R^d}g\right).
\tag{4.5.5}
\end{equation*}

\end{sboxenv}

\textbf{证明}　取有限值可测代表。由柱集可测性，\((u,y)\mapsto f(u)\)
可测。线性变换 \((x,y)\mapsto(x-y,y)\) 可逆，由 §2.4 保持 Lebesgue
可测性，故 \(f(x-y)g(y)\) 联合可测。先对绝对值用 Tonelli，选择先对 \(x\)
积分，使 \(y\) 成为平移参数，得到 \begin{equation*}
\int_{\mathbb R^d}\!\int_{\mathbb R^d}|f(x-y)g(y)|\,dx\,dy
=\int_{\mathbb R^d}|g(y)|\,\|f\|_1\,dy=\|f\|_1\|g\|_1.
\end{equation*} 换回先对 \(y\)
积分的次序，便知卷积几乎处处绝对收敛，且由积分三角不等式得到 (4.5.5)
的估计。现在可对原函数使用 Fubini，先对 \(x\) 积分，平移不变性给出
(4.5.5) 的等式。改变 \(f\) 或 \(g\) 的零测集代表不会改变卷积的 \(L^1\)
等价类，对代表之差应用同一估计即可。

\textbf{例}　取 \(f=g=\mathbf1_{[0,1]}\)。内积分在 \(0\le y\le1\) 且
\(0\le x-y\le1\) 时取单位密度，卷积值就是 \([0,1]\cap[x-1,x]\)
的长度。在 \(0\le x\le1\) 上等于 \(x\)，在 \(1<x\le2\) 上等于
\(2-x\)，其余为零。所得三角形面积为 \(1\)，与 (4.5.5)
中两个区间长度的乘积一致。

\textbf{为何只能保证几乎处处}　令
\(u(t)=|t|^{-1/2}\mathbf1_{\{0<|t|<1\}}\)，在零点取零，则
\(u\in L^1(\mathbb R)\)。但在 \(x=0\) 处，卷积的被积函数为
\(|t|^{-1}\)（\(0<|t|<1\)），故该点的积分为无穷。在例外点补零不影响
(4.5.5)。

\clearpage\section[\S{}4.5.7 练习 3.7 第 1---3
题]{4.5.7 练习 3.7 第 1---3
题}\label{ux7ec3ux4e60-3.7-ux7b2c-13-ux9898}

\Needspace{9\baselineskip}\subsection{问题 1}\label{ux95eeux9898-1}

设 \(f\in L^1([0,1]^2)\)。证明先沿 \(0\le y\le x\) 再对 \(x\)
积分，等于先沿 \(y\le x\le1\) 再对 \(y\) 积分。

\textbf{证明}　令 \(D=\{(x,y):0\le y\le x\le1\}\)。因
\(|f\mathbf1_D|\le|f|\)，Fubini 适用。固定 \(x\in[0,1]\) 得
\(D_x=[0,x]\)，固定 \(y\in[0,1]\) 得 \(D^y=[y,1]\)。于是 \begin{equation*}
\int_0^1\left(\int_0^x f(x,y)\,dy\right)dx
=\int_D f
=\int_0^1\left(\int_y^1 f(x,y)\,dx\right)dy .
\end{equation*}

\Needspace{9\baselineskip}\subsection{问题 2}\label{ux95eeux9898-2}

计算 \(\iint_{[0,\infty)^2}\frac{dx\,dy}{(1+y)(1+x^2y)}\)。

\textbf{解}　函数非负，对每个 \(y>0\) 先用 \(u=x\sqrt y\) 消去分母中
\(x^2\) 的系数。\(y=0\)
的水平截面上内积分无穷，但这个单点不影响外积分。由 Tonelli， \begin{equation*}
\begin{aligned}
\iint_{[0,\infty)^2}\frac{dx\,dy}{(1+y)(1+x^2y)}
&=\int_0^\infty\frac{1}{(1+y)\sqrt y}
       \left(\int_0^\infty\frac{du}{1+u^2}\right)dy\\
&=\frac{\pi}{2}\int_0^\infty\frac{dy}{(1+y)\sqrt y}
=\pi\int_0^\infty\frac{dv}{1+v^2}=\frac{\pi^2}{2}.
\end{aligned}
\end{equation*} 倒数第二个等号用了
\(v=\sqrt y\)。所得有限值也验证了原函数在第一象限可积。

\Needspace{9\baselineskip}\subsection{问题 3}\label{ux95eeux9898-3}

设 \(f\ge0\) 且 \(f\in L^1(0,\infty)\)，令
\(F(x)=\frac1x\int_0^x f(t)\,dt\)。讨论 \(F\) 在 \((0,\infty)\)
上的可积性。

\textbf{解}　教材直接断言 \(F\) 不可积，需补上 \(f\) 不几乎处处为零。若
\(\int_0^\infty f=0\)，非负性给出 \(f=0\) 几乎处处，故对每个 \(x>0\)
均有 \(F(x)=0\)。若总积分为正，由递增收敛，存在 \(a>0\) 使
\(c=\int_0^a f>0\)。这时对 \(x\ge a\)，有 \(F(x)\ge c/x\)，故
\(\int_a^\infty F=\infty\)。两种情形合并，\(F\in L^1\) 当且仅当 \(f=0\)
几乎处处。

也可用换序看见对数从何而来。任取 \(R>a\)，只积区域 \(0<t<a<x<R\)
的非负贡献，由 Tonelli 得
\(\int_a^R F(x)\,dx\ge\int_0^a f(t)\,dt\int_a^R\frac{dx}{x}=c\log(R/a)\)。

\Needspace{10\baselineskip}\section[\S{}4.5.8 练习 3.7 第 4---6
题]{4.5.8 练习 3.7 第 4---6
题}\label{ux7ec3ux4e60-3.7-ux7b2c-46-ux9898}

\Needspace{9\baselineskip}\subsection{问题 4}\label{ux95eeux9898-4}

设 \(E\subset\mathbb R\) 可测，\(f,g\ge0\) 可测且 \(fg\in L^1(E)\)。对
\(t>0\) 令 \(E_t=\{x\in E:g(x)\ge t\}\)、\(F(t)=\int_{E_t}f\)。证明每个
\(F(t)\) 都有限，且 \(\int_0^\infty F(t)\,dt=\int_E fg\)。

\textbf{证明}　在 \(E_t\) 上有 \(tf\le fg\)，故
\(F(t)\le\frac1t\int_E fg<\infty\)。随 \(t\) 增大，\(E_t\) 缩小，所以
\(F\) 单调不增，因而可测。对 \(f(x)\mathbf1_{\{0<t\le g(x)\}}\) 用
Tonelli。先对 \(x\) 积分得 \(F(t)\)，先对 \(t\) 积分则得
\(f(x)g(x)\)，因为相应区间长为
\(g(x)\)。等式得证。这里把严格超水平集换为含端点的水平集，只是给每个竖截面增加一个端点。

\Needspace{9\baselineskip}\subsection{问题 5}\label{ux95eeux9898-5}

设 \(A,B\subset\mathbb R^d\) 可测，\(A-x=\{a-x:a\in A\}\)。证明
\(\int_{\mathbb R^d}m_d((A-x)\cap B)\,dx=m_d(A)m_d(B)\)。

\textbf{证明}　固定 \(x\)，交集中的 \(y\) 同时满足 \(y\in B\) 和
\(x+y\in A\)。所以所求测度是
\(\int\mathbf1_A(x+y)\mathbf1_B(y)\,dy\)。与卷积证明一样，映射
\((x,y)\mapsto(x+y,y)\) 可逆且线性，故被积函数联合可测。先对 \(x\)
积分得 \(\int\mathbf1_A(x+y)\,dx=m_d(A)\)；更精确地说，乘上
\(\mathbf1_B(y)\) 后，内积分在 \(y\in B\) 时为
\(m_d(A)\)，其余为零。Tonelli
给出所求乘积，亦涵盖零测度与无穷测度相乘的情况。

\Needspace{9\baselineskip}\subsection{问题 6}\label{ux95eeux9898-6}

在 \(0\le m(E)<\infty\) 的可测集 \(E\subset\mathbb R\) 上，设 \(f,g\)
为实值可测函数，且 \(H(x,y)=f(x)+g(y)\) 在 \(E\times E\) 上可积。证明
\(f,g\in L^1(E)\)。

\textbf{证明}　若 \(m(E)=0\)，结论立即成立。设 \(0<m(E)<\infty\)。对
\(|H|\) 用 Tonelli，几乎所有 \(y\in E\) 均满足
\(\int_E|f(x)+g(y)|\,dx<\infty\)。因 \(E\) 有正测度，可取一个这样的
\(y_0\)。固定这个截面后，\(g(y_0)\) 是有限常数，故 \begin{equation*}
\int_E|f(x)|\,dx
\le \int_E|f(x)+g(y_0)|\,dx+m(E)|g(y_0)|<\infty .
\end{equation*} 交换 \(x,y\) 的角色，即得
\(g\in L^1(E)\)。有限测度用在常数项的积分上；若 \(E=\mathbb R\)，取
\(f=1\)、\(g=-1\)，则 \(H=0\) 可积而两个常函数均不可积。

\end{document}
